% =====================================================================
%  FICHE 11 — THEME 2 — Exercices MOYEN — Oxydant, réducteur, demi-équations
% =====================================================================
\documentclass[11pt,a4paper]{article}
\usepackage[utf8]{inputenc}
\usepackage[T1]{fontenc}
\usepackage{lmodern}
\usepackage[french]{babel}
\usepackage{amsmath,amssymb}
\usepackage[version=4]{mhchem}
\usepackage{siunitx}
\sisetup{output-decimal-marker={,},per-mode=symbol}
\usepackage{xcolor}
\usepackage{array}
\usepackage{booktabs}
\usepackage{enumitem}
\usepackage[most]{tcolorbox}
\usepackage[a4paper,margin=2.1cm,headheight=26pt,headsep=10pt]{geometry}
\usepackage{fancyhdr}
\usepackage[hidelinks]{hyperref}

\definecolor{primary}{RGB}{146,95,160}
\definecolor{primarydark}{RGB}{92,52,104}
\newcommand{\NO}{\ensuremath{\mathrm{N.O.}}}

\newcounter{exo}
\newtcolorbox{exo}[1][]{enhanced,breakable,boxrule=0.9pt,arc=2pt,colback=primary!4,
  colframe=primary,fonttitle=\bfseries,coltitle=white,left=3mm,right=3mm,top=2.2mm,bottom=2.2mm,
  title={Exercice~\refstepcounter{exo}\theexo\ifx\relax#1\relax\relax\else\ \textmd{--- \itshape #1}\fi}}

\pagestyle{fancy}\fancyhf{}
\fancyhead[L]{\small\textcolor{primarydark}{\textbf{Fiche 11 · Thème 2} — Oxydant, réducteur, demi-équations}}
\fancyhead[R]{\small\textcolor{primarydark}{\textsc{Moyen}}}
\fancyfoot[C]{\small\thepage}
\renewcommand{\headrulewidth}{0.8pt}
\renewcommand{\headrule}{\hbox to\headwidth{\color{primary}\leaders\hrule height \headrulewidth\hfill}}

\begin{document}
\begin{center}
{\LARGE\bfseries\color{primarydark}Thème 2 — Niveau Moyen}\\[2pt]
{\small\itshape Compter les électrons, reconnaître le redox, repérer oxydant, réducteur et spectateurs.}\\[-4pt]
{\color{primary}\rule{\textwidth}{1pt}}
\end{center}

\begin{exo}[compter les électrons de trois couples]
Pour chaque couple, calcule le nombre d'électrons échangés à partir de la variation de \NO, puis
écris la demi-équation de réduction (sans encore équilibrer O et H) :
\begin{enumerate}[label=(\alph*),leftmargin=2em,topsep=1pt,itemsep=1pt]
  \item \ce{NO3^-}/\ce{NO} ;
  \item \ce{MnO4^-}/\ce{Mn^{2+}} ;
  \item \ce{Cr2O7^{2-}}/\ce{Cr^{3+}} \quad\emph{(attention au nombre d'atomes de chrome !)}.
\end{enumerate}
\end{exo}

\begin{exo}[trier les réactions redox — d'après annales 2024]
Parmi les cinq réactions suivantes, indique lesquelles sont des réactions d'\textbf{oxydo-réduction}.
Justifie en repérant un élément dont le \NO\ change (ou en montrant qu'aucun ne change).
\begin{enumerate}[label=\arabic*),leftmargin=2.2em,topsep=1pt,itemsep=1pt]
  \item \ce{NaCl + AgNO3 <=> AgCl + NaNO3}
  \item \ce{NaHSO3 + Na2CO3 <=> NaHCO3 + Na2SO3}
  \item \ce{3 CoCl3 + MnO2 + 2 H2O + KCl <=> 3 CoCl2 + KMnO4 + 4 HCl}
  \item \ce{2 Na3PO4 + 3 BaCl2 <=> Ba3(PO4)2 + 6 NaCl}
  \item \ce{H2O2 + MnCl2 <=> MnO2 + 2 HCl}
\end{enumerate}
\end{exo}

\begin{exo}[oxydant et réducteur dans une équation globale — d'après livre QCM 9]
On donne la réaction : $\ce{4 HCl + MnO2 -> MnCl2 + Cl2 + 2 H2O}$.
\begin{enumerate}[label=(\alph*),leftmargin=2em,topsep=1pt,itemsep=1pt]
  \item Suis le \NO\ du manganèse et celui du chlore entre réactifs et produits.
  \item Quel réactif est l'\textbf{oxydant} ? Lequel est le \textbf{réducteur} ?
\end{enumerate}
\end{exo}

\begin{exo}[ne pas oublier de multiplier les électrons]
Détermine le nombre total d'électrons échangés dans chacune des demi-équations de réduction
suivantes (pense au nombre d'atomes concernés) :
\[
\ce{I2}/\ce{I^-},\qquad \ce{O2}/\ce{H2O},\qquad \ce{S4O6^{2-}}/\ce{S2O3^{2-}}.
\]
\end{exo}

\begin{exo}[les ions spectateurs — d'après livre QCM 10]
On donne la réaction en solution :
\[
\ce{NaNO3 + 3 FeCl2 + 4 HCl -> NO + 3 FeCl3 + 2 H2O + NaCl}.
\]
\begin{enumerate}[label=(\alph*),leftmargin=2em,topsep=1pt,itemsep=1pt]
  \item Dresse la liste de tous les ions présents (réactifs et produits).
  \item Identifie les deux couples redox mis en jeu (éléments dont le \NO\ varie).
  \item Quels sont les \textbf{ions spectateurs} ?
\end{enumerate}
\end{exo}

\end{document}
